
%% Use the "normalphoto" option if you want a normal photo instead of cropped to a circle
% \documentclass[10pt,a4paper,normalphoto]{altacv}

\documentclass[10pt,a4paper,ragged2e,withhyper]{altacv}
%% AltaCV uses the fontawesome5 and simpleicons packages.
%% See http://texdoc.net/pkg/fontawesome5 and http://texdoc.net/pkg/simpleicons for full list of symbols.

% Change the page layout if you need to
%\geometry{left=1.25cm,right=1.25cm,top=1.5cm,bottom=1.5cm,columnsep=1.2cm}
\geometry{left=1cm,right=1cm,top=1cm,bottom=1cm,columnsep=1cm}
% The paracol package lets you typeset columns of text in parallel
\usepackage{paracol}

\usepackage{ragged2e}

% xelatex -shell-escape -output-driver="xdvipdfmx -z 0" sample.tex
\ifxetexorluatex
  % If using xelatex or lualatex:
  \setmainfont{Roboto Slab}
  \setsansfont{Lato}
  \renewcommand{\familydefault}{\sfdefault}
\else
  % If using pdflatex:
  \usepackage[rm]{roboto}
  \usepackage[defaultsans]{lato}
  % \usepackage{sourcesanspro}
  \renewcommand{\familydefault}{\sfdefault}
\fi

% Change the colours if you want to
\definecolor{SlateGrey}{HTML}{2E2E2E}
\definecolor{LightGrey}{HTML}{666666}
\definecolor{DarkPastelRed}{HTML}{450808}
\definecolor{PastelRed}{HTML}{8F0D0D}
\definecolor{GoldenEarth}{HTML}{E7D192}
\colorlet{name}{black}
\colorlet{tagline}{PastelRed}
\colorlet{heading}{DarkPastelRed}
\colorlet{headingrule}{GoldenEarth}
\colorlet{subheading}{PastelRed}
\colorlet{accent}{PastelRed}
\colorlet{emphasis}{SlateGrey}
\colorlet{body}{LightGrey}

% Change some fonts, if necessary
\renewcommand{\namefont}{\Huge\rmfamily\bfseries}
\renewcommand{\personalinfofont}{\footnotesize}
\renewcommand{\cvsectionfont}{\LARGE\rmfamily\bfseries}
\renewcommand{\cvsubsectionfont}{\large\bfseries}

% Ajoute ça avant \begin{document}
\newcommand{\kw}[1]{\textbf{\color{name}#1}}

% Change the bullets for itemize and rating marker
% for \cvskill if you want to
\renewcommand{\cvItemMarker}{{\small\textbullet}}
\renewcommand{\cvRatingMarker}{\faCircle}
% ...and the markers for the date/location for \cvevent
% \renewcommand{\cvDateMarker}{\faCalendar*[regular]}
% \renewcommand{\cvLocationMarker}{\faMapMarker*}

\begin{document}
\name{Lucas Noirot (Petit)}

\tagline{Ingénieur R\&D en Algorithmique \& Mathématiques Appliquées}
%% You can add multiple photos on the left or right
%\photoR{4cm}{img_7056}
%\photoL{2.8cm}{hovercode}

\personalinfo{%
  \email{lucas.noirot-petit@proton.me}
  \phone{(+33) 685263667}
  \location{Montpellier, France}
  \homepage{enpeluche.github.io}
  \linkedin{enpeluche}
  \github{enpeluche}
  %% You can add your own arbitrary detail with
  %% \printinfo{symbol}{detail}[optional hyperlink prefix]
  % \printinfo{\faPaw}{Hey ho!}[https://example.com/]

  %% Or you can declare your own field with
  %% \NewInfoFiled{fieldname}{symbol}[optional hyperlink prefix] and use it:
  % \NewInfoField{gitlab}{\faGitlab}[https://gitlab.com/]
  % \gitlab{your_id}
  %%
  %% For services and platforms like Mastodon where there isn't a
  %% straightforward relation between the user ID/nickname and the hyperlink,
  %% you can use \printinfo directly e.g.
  % \printinfo{\faMastodon}{@username@instace}[https://instance.url/@username]
  %% But if you absolutely want to create new dedicated info fields for
  %% such platforms, then use \NewInfoField* with a star:
  % \NewInfoField*{mastodon}{\faMastodon}
  %% then you can use \mastodon, with TWO arguments where the 2nd argument is
  %% the full hyperlink.
}


\columnratio{0.35}

\begin{paracol}{2}

  %%%%%%%%%%%%%%%%%%%%%%%%%%%%%%%%%%%%%%%%%%%%%%%%%%%%%%%%%%%%%%%%%
  \makecvheaderright

  \cvsection{À propos de moi}

  \begin{quote}
    Passionné par l'informatique et les mathématiques, je trouve mon équilibre dans l'abstraction, où
    les idées prennent forme et deviennent concrètes.
  \end{quote}


  \cvsection{Compétences}

  \cvsubsection{Sûreté \& Cryptographie}
  \cvtag{Méthodes formelles (Coq)}
  \cvtag{Preuve de sécurité}

  \vspace{0.4em}

  \cvsubsection{Optimisation \& Modélisation}
  \cvtag{PLNE (MIP)}
  \cvtag{Programmation par contraintes}
  \cvtag{Heuristiques}
  \cvtag{Analyse de complexité}

  \vspace{0.4em}

  \cvsubsection{Développement}
  \cvtag{C/C++} \hfill
  \cvtag{Rust} \hfill
  \cvtag{OCaml} \hfill
  \cvtag{Python} \\
  \cvtag{SageMath} \hfill
  \cvtag{Java} \hfill
  \cvtag{HTML/JS}

  \vspace{0.4em}

  \cvsubsection{Outils \& Environnement}
  \cvtag{Linux} \hfill
  \cvtag{Git / GitHub} \hfill
  \cvtag{Docker} \\
  \cvtag{\LaTeX} \hfill
  \cvtag{Visualisation}
  %%%%%%%%%%%%%%%%%%%%%%%%%%%%%%%%%%%%%%%%%%%%%%%%%%%%%%%%%%%%%%%%%

  \cvsection{Langues}

  \begin{itemize}[label=\textcolor{accent}{\tiny\faCircle}]
    \item \textbf{Français} (Natif)
    \item \textbf{Anglais} (C1 -- Technique)
  \end{itemize}

  \cvsection{Autre}

  \cvsubsection{Qualités personnelles}
  \cvtag{Rigoureux} \hfill
  \cvtag{Abstraction} \hfill
  \cvtag{Autodidacte}\\

  \cvtag{Pro-actif} \hfill
  \cvtag{Esprit d'équipe} \hfill
  \cvtag{Créatif}

  \vspace{0.4em}

  \cvsubsection{Centres d'intérêt}
  \cvtag{Origami}
  \cvtag{Jeux de société}
  \cvtag{Astronomie}

  %%%%%%%%%%%%%%%%%%%%%%%%%%%%%%%%%%%%%%%%%%%%%%%%%%%%%%%%%%%%%%%%%
  %
  \switchcolumn
  %
  %%%%%%%%%%%%%%%%%%%%%%%%%%%%%%%%%%%%%%%%%%%%%%%%%%%%%%%%%%%%%%%%%

  %%%%%%%%%%%%%%%%%%%%%%%%%%%%%%%%%%%%%%%%%%%%%%%%%%%%%%%%%%%%%%%%%
  \makecvheaderleft

  %%%%%%%%%%%%%%%%%%%%%%%%%%%%%%%%%%%%%%%%%%%%%%%%%%%%%%%%%%%%%%%%%
  \cvsection{Formations}

  \cveventt{Master en Informatique - Algorithmique}
  {Université de Montpellier}
  {\smash{\parbox[t]{3em}{\centering
        2025
        \\[3pt]
        \vrule height 27pt width 0.5pt}}}
  {Complexité -- Cryptographie -- Calcul formel -- Optimisation}

  \cveventt{Master 1 en Mathématiques}
  {Université de Bourgogne}
  {\smash{\parbox[t]{3em}{\centering
        2023
        \\[3pt]
        \vrule height 27pt width 0.5pt}}}
  {Algèbre linéaire -- Théorie des corps -- Probabilités }

  \cveventt{Licence en Mathématiques}
  {Université de Bourgogne}
  {\smash{\parbox[t]{3em}{\centering
        2022
        \\[3pt]
        \vrule height 11pt width 0.5pt \\[-1pt] 2019}}}
  {Classé 5\textsuperscript{e} sur 55}

  %%%%%%%%%%%%%%%%%%%%%%%%%%%%%%%%%%%%%%%%%%%%%%%%%%%%%%%%%%%%%%%%%
  \cvsection{Expérience}

  \experiences{Stage de Recherche en Cryptographie}
  {Laboratoire du LIRMM -- Équipe ECO}{
  \smash{\parbox[t]{3em}{\centering
  2025
  \\[3pt]
  {\scriptsize\itshape (5 mois)}}}}
  {\vspace{0.05cm}

  Réduction de réseaux : Adaptations d’idées provenant du cas polynomial.

  Développement d'algorithmes de réduction sous \kw{SageMath}. Analyse de complexité et comparaison des performances entre arithmétique polynomiale et entière.

  {\small \faGithub \color{accent}  \href{https://enpeluche.github.io/lattice-reduction-internship/}{enpeluche.github.io/lattice-reduction-internship/}}
  }

  \vspace{0.4cm}

  \experiences{Professeur particulier}
  {Dijon, freelance}
  {2022}
  {Mathématiques et physique niveau collège et lycée}

  %%%%%%%%%%%%%%%%%%%%%%%%%%%%%%%%%%%%%%%%%%%%%%%%%%%%%%%%%%%%%%%%%
  \cvsection{Projets personnels}

  \project{Optimisation de Tournées \& Planification (WSRP)}
  {Conception d'un solveur hybride pour le routage de techniciens sous contraintes de compétences. Implémentation d'une \kw{heuristique} en \kw{Java} couplée à une programmation par contraintes (\kw{OR-Tools}).}
  {enpeluche/wsrp-lns-solver}

  \project{Coq-Verified Red-Black Tree}
  {Formalisation et preuve mathématique des invariants d'équilibrage dans \kw{Coq}, avec extraction automatique vers un module \kw{OCaml} performant et certifié sans bug.}
  {enpeluche/coq-verified-rbt}

  \project{Architecture de Moteur de Résolution de Contraintes (CSP)}
  {Développement \textit{from scratch} en \kw{C++} d'un moteur de propagation. Implémentation des algorithmes d'Arc-Consistance (\kw{AC-3}).}
  {enpeluche/constraint-puzzle-solver}

  \project{Protocole Cryptographique ZKP (Zero-Knowledge Proof)}
  {Implémentation du protocole d'authentification de \kw{Schnorr sur courbes elliptiques}. Prototypage algébrique sous \kw{SageMath} et développement d'une version performante en \kw{C++}.}
  {enpeluche/zkp-schnorr-cpp}

  \project{Messagerie Post-Quantique (Réseaux Euclidiens)}
  {Implémentation d'un chat chiffré résistant aux ordinateurs quantiques. Développement \kw{C++} d'un protocole d'échange de clés basé sur le problème \kw{LWE (Learning With Errors)} et les réseaux euclidiens. Architecture Client/Serveur asynchrone.}
  {https://github.com/enpeluche/}

\end{paracol}

\end{document}